% \documentclass[aspectratio=169,notes]{beamer}
\documentclass[aspectratio=169]{beamer}
\usetheme[faculty=phil]{fibeamer}
\usepackage{polyglossia}
\setmainlanguage{english} %% main locale instead of `english`, you
%% can typeset the presentation in either Czech or Slovak,
%% respectively.
\setotherlanguages{russian} %% The additional keys allow
%%
%%   \begin{otherlanguage}{czech}   ... \end{otherlanguage}
%%   \begin{otherlanguage}{slovak}  ... \end{otherlanguage}
%%
%% These macros specify information about the presentation
\title[Artificial Intelligence Software (AIS)]{Artificial Intelligence Software, Lecture 8} %% that will be typeset on the
\subtitle{Experiments \& Hyperparameters
\\ \  \\ \  } %% title page.
\author{Oleg Bulichev}
%% These additional packages are used within the document:
\usepackage{ragged2e}  % `\justifying` text
\usepackage{booktabs}  % Tables
\usepackage{tabularx}
\usepackage{tikz}      % Diagrams
\usetikzlibrary{calc, shapes, backgrounds}
\usetikzlibrary{decorations.pathreplacing,calligraphy,calc,graphs}
\usepackage{amsmath, amssymb}
\usepackage{url}       % `\url`s
\usepackage{listings}  % Code listings
\usepackage{subfigure}
% \usepackage{floatrow}
% \usepackage{subcaption}
\usepackage{mathtools}
\usepackage{todonotes}
\usepackage{fontspec}
\usepackage{multicol}
\usepackage{pdfpages}
\usepackage{wrapfig}
\usepackage{qrcode}
\usepackage{animate}
\usepackage{booktabs}
\usepackage{multirow}
\usepackage{graphicx}
\usepackage{colortbl}

\graphicspath{{resources/}}
\frenchspacing

\setbeamertemplate{caption}{\raggedright\insertcaption\par}
\usetikzlibrary{graphs}
\usetikzlibrary{positioning,arrows.meta}

% \usepackage[backend=biber,style=ieee,autocite=footnote]{biblatex}
% \addbibresource{biblio.bib}
% \DefineBibliographyStrings{english}{%
%   bibliography = {References},}

\newcommand{\oleg}[2][] {\todo[color=red, #1] {OLEG:\\ #2}}
\newcommand{\fbckg}[1]{\usebackgroundtemplate{\includegraphics[width=\paperwidth]{#1}}}%frame background

\usepackage[framemethod=TikZ]{mdframed}
\newcommand{\dbox}[1]{
\begin{mdframed}[roundcorner=3pt, backgroundcolor=yellow, linewidth=0]
\vspace{1mm}
{#1}
\vspace{1mm}
\end{mdframed}
}

\newcommand{\imgmock}[2]{
\fbox{\parbox[c][0.48\textheight][c]{0.92\linewidth}{\centering\textbf{Mock image placeholder}\\[2mm]#1\\[2mm]\footnotesize #2}}
}

\begin{document}
\setlength{\abovedisplayskip}{0pt}
\setlength{\belowdisplayskip}{0pt}
\setlength{\abovedisplayshortskip}{0pt}
\setlength{\belowdisplayshortskip}{0pt}

\fbckg{fibeamer/figs/title_page.png}
\frame[c]{\setcounter{framenumber}{0}
    \usebeamerfont{title}%
    \usebeamercolor[fg]{title}%
    \begin{minipage}[b][6.5\baselineskip][b]{\textwidth}%
        \textcolor{black}{\raggedright\inserttitle}
    \end{minipage}
    % \vskip-1.5\baselineskip

    \usebeamerfont{subtitle}%
    \usebeamercolor[fg]{framesubtitle}%
    \begin{minipage}[b][3\baselineskip][b]{\textwidth}
        \raggedright%
        \insertsubtitle%
    \end{minipage}
    \vskip.25\baselineskip
}
%   \frame[c]{\maketitle}

\fbckg{fibeamer/figs/common.png}

\note{\scriptsize \begin{itemize}
        \item \
    \end{itemize}}


\begin{frame}[t]{Lecture Plan}
\begin{enumerate}
    \item Motivation: why tracking experiments matters
    \item Experiment tracking with \textbf{ClearML}
    \item Visualizing training with \textbf{TensorBoard}
    \item What are hyperparameters and search strategies
    \item Automated HPO with \textbf{Optuna}
    \item Demo: YOLO + ClearML + Optuna + TensorBoard\\
    {\footnotesize\texttt{demo/notebooks/demo\_hparams.ipynb}}
\end{enumerate}
\end{frame}

\begin{frame}[t]{The Experiment Problem}
A typical ML workflow involves running many experiments with different configurations.
\vspace{2mm}
\textbf{Without tracking:}
\begin{itemize}
    \item You cannot reproduce a result from last week.
    \item You lose best checkpoints and their exact configurations.
    \item Comparing two runs objectively is impossible.
\end{itemize}
\vspace{2mm}
\textbf{Solution:} experiment tracking systems that log everything automatically ---
code version, packages, config, metrics, artifacts.
\end{frame}

\begin{frame}[t,fragile]{ClearML: Task Initialization}
\begin{lstlisting}[language=Python, basicstyle=\footnotesize\ttfamily]
from clearml import Task

task = Task.init(
    project_name="YOLO-Roboflow",   # groups tasks on dashboard
    task_name="trial-lr0.001",      # human-readable run name
    task_type=Task.TaskTypes.training,  # optional: training/testing/...
    reuse_last_task_id=False        # always create a new task
)
logger = task.get_logger()
\end{lstlisting}
\vspace{1mm}
Auto-captured without any extra code: git diff, Python packages, stdout/stderr, hostname, GPU/CPU info, script source.\\
\textbf{Rule:} call \texttt{Task.init()} before any other ML code.
\end{frame}

\begin{frame}[t,fragile]{ClearML: Logging Scalars}
\texttt{logger.report\_scalar(title, series, value, iteration)}
\begin{lstlisting}[language=Python, basicstyle=\footnotesize\ttfamily]
# title   -- chart name on dashboard (groups related series)
# series  -- line name within that chart
# iteration -- x-axis value (step / epoch)

logger.report_scalar("Loss",     "train", value=0.42, iteration=5)
logger.report_scalar("Loss",     "val",   value=0.55, iteration=5)
logger.report_scalar("mAP50",    "val",   value=0.71, iteration=5)
\end{lstlisting}
\vspace{1mm}
Result: a ``Loss'' chart with two lines (\textit{train} and \textit{val}) and a separate ``mAP50'' chart --- all on the same dashboard page.
\end{frame}

\begin{frame}[t,fragile]{ClearML: Logging Images \& Plots}
\begin{lstlisting}[language=Python, basicstyle=\footnotesize\ttfamily]
logger.report_image(
    title="Predictions", series="batch_0",
    iteration=epoch, image=pred_img)
# matplotlib figure
import matplotlib.pyplot as plt
fig, ax = plt.subplots()
ax.plot(losses)
logger.report_matplotlib_figure(
    title="Loss curve", series="train",
    iteration=epoch, figure=fig)
plt.close(fig)
\end{lstlisting}
Images appear under Task $\rightarrow$ \textit{Debug Samples}; figures under \textit{Plots}.
\end{frame}

\begin{frame}[t,fragile]{ClearML: Hyperparameters \& Artifacts}
\vspace*{-0.7cm}
\begin{lstlisting}[language=Python, basicstyle=\footnotesize\ttfamily]
# --- Hyperparameters (Task > Configuration) ---
hparams = {"lr0": 1e-3, "batch": 16, "epochs": 30}
task.connect(hparams, name="training_config")
# dict is mutable: ClearML may override values
# from the dashboard before training starts

# --- Upload any file as an artifact ---
task.upload_artifact("confusion_matrix", "cm.png")
task.upload_artifact("val_predictions",  predictions_df)

# --- Register a trained model ---
from clearml import OutputModel
output_model = OutputModel(task=task, framework="PyTorch")
output_model.update_weights("best.pt", auto_delete_file=False)
\end{lstlisting}
Artifacts and models are versioned and downloadable from the dashboard.
\end{frame}

\begin{frame}[t,fragile]{ClearML: Tracking Hyperparameters (demo pattern)}
\begin{lstlisting}[language=Python, basicstyle=\footnotesize\ttfamily]
hparams = {"epochs": 5, "batch": 16, "lr0": 1e-3}

# Appears under Task > Configuration on the dashboard
task.connect(hparams, name="training_config")
\end{lstlisting}

\vspace{2mm}
In the demo, this is called per Optuna trial:
\begin{lstlisting}[language=Python, basicstyle=\footnotesize\ttfamily]
task.connect(yolo_hparams,
             name=f"hparams_trial_{trial.number}")
\end{lstlisting}
Multiple calls with different \texttt{name} create separate config sections on the dashboard.
\end{frame}


\begin{frame}[t]{What is TensorBoard?}
TensorBoard is a web-based visualization tool for monitoring training runs.

\vspace{2mm}
Key views:
\begin{columns}[T,onlytextwidth]
    \begin{column}{0.49\textwidth}
        \begin{itemize}
            \item \textbf{Scalars}: loss and metric curves
            \item \textbf{Images}: samples, augmentations, predictions
            \item \textbf{Histograms}: weight and gradient distributions
        \end{itemize}
    \end{column}
    \begin{column}{0.49\textwidth}
        \begin{itemize}
            \item \textbf{Graph}: model computational graph
            \item \textbf{Projector}: embedding visualization (t-SNE)
            \item \textbf{HParams}: compare HPO trials side-by-side
        \end{itemize}
    \end{column}
\end{columns}
\vspace{2mm}
Launch: \texttt{tensorboard --logdir runs/}
\end{frame}

\begin{frame}[t,fragile]{TensorBoard: Writing Logs}
\vspace*{-0.7cm}
\begin{lstlisting}[language=Python, basicstyle=\footnotesize\ttfamily]
from torch.utils.tensorboard import SummaryWriter

writer = SummaryWriter("runs/cnn_training")

writer.add_scalar("Loss/train",   loss.item(), epoch)
writer.add_scalar("Accuracy/val", val_acc,     epoch)
writer.add_graph(model, dummy_input)

for name, param in model.named_parameters():
    writer.add_histogram(name, param, epoch)

writer.close()
\end{lstlisting}
\vspace{1mm}
See \texttt{demo/notebooks/tensorboard\_showcase/showcase.ipynb}\\for a working CNN example with all TensorBoard views enabled.
\end{frame}

\begin{frame}[t]{What Are Hyperparameters?}
Hyperparameters are \textit{values set before training} that are not learned from data.

\vspace{2mm}
\begin{tabular}{@{}p{3.0cm}p{8.2cm}@{}}
\toprule
Group & Examples \\
\midrule
Optimization  & Learning rate, weight decay, optimizer type \\
Batching      & Batch size, gradient accumulation steps \\
Architecture  & Dropout rate, number of layers, activation functions \\
Regularization & Augmentation strength, label smoothing \\
Schedule      & Warmup epochs, total epochs, early-stop patience \\
\bottomrule
\end{tabular}

Wrong hyperparameters can prevent convergence entirely --- even with a correct architecture.
\end{frame}

\begin{frame}[t]{Hyperparameter Search Strategies}
\begin{tabular}{@{}p{2.8cm}p{8.5cm}@{}}
\toprule
Strategy & Description \\
\midrule
\textbf{Grid Search}   & Try all combinations. Exhaustive; scales exponentially with dimensions. \\
\textbf{Random Search} & Sample uniformly at random. Surprisingly competitive for high-dimensional spaces. \\
\textbf{Bayesian (TPE)} & Build a probabilistic model of $f(\text{hparams})$; suggest promising regions. Most efficient. \\
\textbf{Pruning (ASHA)} & Start many trials cheaply; kill under-performers early. Saves 3--10$\times$ compute. \\
\bottomrule
\end{tabular}

Practical default: \textbf{Bayesian search + ASHA pruning}.
\end{frame}

\begin{frame}[t,fragile]{Optuna: Core Concepts}
\vspace*{-0.6cm}
\begin{lstlisting}[language=Python, basicstyle=\footnotesize\ttfamily]
import optuna

def objective(trial):
    lr    = trial.suggest_float("lr", 1e-5, 1e-2, log=True)
    batch = trial.suggest_categorical("batch", [8, 16, 32])

    score = train_and_evaluate(lr, batch)  # returns mAP / accuracy
    return score

study = optuna.create_study(direction="maximize")
study.optimize(objective, n_trials=20)

print(study.best_params)  # {'lr': 0.0031, 'batch': 16}
\end{lstlisting}
\vspace{1mm}
\textbf{Study}: optimization run (history of all trials). \quad \textbf{Trial}: one \texttt{objective} call with a sampled config.
\end{frame}

\begin{frame}[t,fragile]{Optuna: Search Space API}
\vspace*{-0.6cm}
\begin{lstlisting}[language=Python, basicstyle=\footnotesize\ttfamily]
# Log-uniform float -- ideal for learning rates
lr = trial.suggest_float("lr", 1e-5, 1e-2, log=True)

# Uniform float -- dropout, weight decay, etc.
dropout = trial.suggest_float("dropout", 0.0, 0.5)

# Integer -- number of layers, warmup steps
n_layers = trial.suggest_int("n_layers", 1, 5)

# Categorical -- discrete choices
optimizer_name = trial.suggest_categorical(
    "optimizer", ["adam", "sgd", "adamw"]
)
\end{lstlisting}
In the demo: \texttt{lr} is log-uniform, \texttt{batch} $\in\{8, 16\}$ is categorical.
\end{frame}

\begin{frame}[t]{Optuna: Pruning (Early Stopping of Bad Trials)}
Pruning monitors intermediate values and terminates weak trials early.

\vspace{2mm}
\textbf{Workflow:}
\begin{enumerate}
    \item Report intermediate value during training:\\
    \texttt{trial.report(mAP50, step=epoch)}
    \item Check whether to stop:\\
    \texttt{if trial.should\_prune(): raise optuna.TrialPruned()}
    \item Optuna's pruner decides based on all trials seen so far.
\end{enumerate}

\vspace{2mm}
\textbf{ASHA} (Asynchronous Successive Halving) --- standard pruner: promotes top-$k$\% of trials at each stage; discards the rest.
\vspace{1mm}

Result: spend GPU time only on configurations that are promising early on.
\end{frame}

\begin{frame}[t,fragile]{Optuna + TensorBoard Integration}
\vspace*{-0.6cm}
\begin{lstlisting}[language=Python, basicstyle=\footnotesize\ttfamily]
from optuna.integration.tensorboard import TensorBoardCallback

tb_callback = TensorBoardCallback(
    "runs/tb_logs",
    metric_name="mAP50")

study.optimize(
    objective,
    n_trials=3,
    callbacks=[tb_callback])
\end{lstlisting}
\vspace{1mm}
Each trial is logged as a separate run under \texttt{runs/tb\_logs/}.\\Open TensorBoard $\rightarrow$ \textit{HParams} tab to compare all trials side-by-side.
\end{frame}

\begin{frame}[t]{Demo Pipeline Architecture}
\centering
\begin{tikzpicture}[
    box/.style={draw, rounded corners, align=center,
                minimum width=2.6cm, minimum height=0.8cm, fill=blue!8, font=\small},
    cbox/.style={draw, rounded corners, align=center,
                minimum width=2.6cm, minimum height=0.8cm, fill=green!12, font=\small},
    obox/.style={draw, rounded corners, align=center,
                minimum width=2.6cm, minimum height=0.8cm, fill=orange!14, font=\small},
    arr/.style={-{Latex[length=2.5mm]}, thick}
]
\node[box]  (dataset) {Dataset\\Valentines Chocolates};
\node[cbox, right=0.6cm of dataset]  (task)   {ClearML\\Task.init()};
\node[obox, right=0.6cm of task]     (optuna) {Optuna\\study.optimize()};
\node[box,  below=0.75cm of optuna]  (yolo)   {YOLO training\\(per trial, 5 ep.)};
\node[cbox, below=0.75cm of task]    (log)    {ClearML\\task.connect()\\report\_scalar};
\node[obox, below=0.75cm of dataset] (tb)     {TensorBoard\\TensorBoardCallback};
\node[box,  right=0.6cm of yolo]     (best)   {Best params\\final retrain 30 ep.};

\draw[arr] (dataset) -- (task);
\draw[arr] (task)    -- (optuna);
\draw[arr] (optuna)  -- (yolo);
\draw[arr] (yolo)    -- (log);
\draw[arr] (yolo)    -- (tb);
\draw[arr] (yolo)    -- (best);
\end{tikzpicture}
\end{frame}

\begin{frame}[t]{Demo Walkthrough: \texttt{demo\_hparams.ipynb}}
File: \texttt{demo/notebooks/demo\_hparams.ipynb}
\vspace{1mm}
\begin{enumerate}
    \item \textbf{Init ClearML Task} --- \texttt{Task.init("YOLO-Roboflow", "Valentines-Chocolates")}
    \item \textbf{Load dataset} --- 22-class Valentines Chocolates in YOLOv8 format
    \item \textbf{Define Optuna objective} --- search \texttt{lr} ($10^{-5}$--$10^{-2}$, log-uniform) and \texttt{batch} ($\in\{8,16\}$)
    \item \textbf{Per-trial logging} --- \texttt{task.connect(hparams)} + \texttt{logger.report\_scalar("optuna", "mAP50", ...)}
    \item \textbf{TensorBoard callback} --- writes each trial to \texttt{runs/tb\_logs/trial-N/}
    \item \textbf{Final training} --- 30 epochs with \texttt{study.best\_params}
    \item \textbf{Inference + logging} --- predictions annotated and logged to ClearML
\end{enumerate}
\vspace{1mm}
Run: \texttt{cd demo \&\& pixi run jupyter notebook}
\end{frame}

\begin{frame}[t]{Practical HPO Workflow}
\begin{enumerate}
    \item \textbf{Baseline first}: train with default hyperparameters; log everything to ClearML.
    \item \textbf{Identify key hyperparameters}: 5--10 random trials reveal which params matter most (usually \texttt{lr} and \texttt{batch}).
    \item \textbf{Narrow and refine}: run Bayesian search (TPE) in the promising region.
    \item \textbf{Enable pruning}: ASHA eliminates wasted compute in the first epochs of bad trials.
    \item \textbf{Compare in ClearML}: dashboard side-by-side comparison across all trials.
    \item \textbf{Final run}: retrain with the best config for the full epoch budget.
\end{enumerate}
\end{frame}

\begin{frame}[t]{Summary}
\begin{itemize}
    \item Experiment tracking is not optional --- it is the foundation of reproducible ML.
    \item \textbf{ClearML} auto-captures code, config, metrics, and artifacts for every run.
    \item \textbf{TensorBoard} gives live visual monitoring: loss curves, weight histograms, embeddings.
    \item \textbf{Optuna} turns hyperparameter search into systematic Bayesian optimization.
    \item Pruning (ASHA) makes HPO practical: spend GPU time only on promising trials.
    \item The demo (\texttt{demo\_hparams.ipynb}) shows the complete pipeline on a real detection dataset.
\end{itemize}
\end{frame}



\fbckg{fibeamer/figs/last_page.png}
\frame[plain]{}

\end{document}
